\documentclass[10pt,a4paper]{amsart}
\usepackage{amsmath,amssymb,amsthm,graphicx}
% Retained preamble for the Stacks Project chapter {CH} (shared project preamble).
% The chapter file pinned-original-source.tex only carries \input{{preamble}} on its first line,
% so no preamble text lives inside the chapter itself.  The Stacks Project uses one single shared
% preamble for every chapter: the snapshot stages it three times (sources/stacks-chow/preamble.tex,
% sources/stacks-etale-cohomology/preamble.tex and focused-seeds/AGP/sources/stacks-snapshot-20261011/
% preamble.tex) and those copies are identical in size.  The block below is copied mechanically,
% line for line, from sources/stacks-chow/preamble.tex lines 14-18 and 154-217: the Xy-pic setup and
% the native theorem declarations together with the mathematical macros.  Nothing was added, removed
% or redefined; the wrapper class and the amsmath/amssymb/amsthm/graphicx base packages are supplied
% by the assembler.
% For commutative diagrams we use Xy-pic
% Chapter: stacks-cohomology.tex -- Cohomology of Sheaves.
\usepackage[all]{xy}

% We use 2cell for 2-commutative diagrams.
\xyoption{2cell}

% Theorem environments.
%
\theoremstyle{plain}
\newtheorem{theorem}[subsection]{Theorem}
\newtheorem{proposition}[subsection]{Proposition}
\newtheorem{lemma}[subsection]{Lemma}

\theoremstyle{definition}
\newtheorem{definition}[subsection]{Definition}
\newtheorem{example}[subsection]{Example}
\newtheorem{exercise}[subsection]{Exercise}
\newtheorem{situation}[subsection]{Situation}

\theoremstyle{remark}
\newtheorem{remark}[subsection]{Remark}
\newtheorem{remarks}[subsection]{Remarks}

\numberwithin{equation}{subsection}

% Macros
%
\def\lim{\mathop{\mathrm{lim}}\nolimits}
\def\colim{\mathop{\mathrm{colim}}\nolimits}
\def\Spec{\mathop{\mathrm{Spec}}}
\def\Hom{\mathop{\mathrm{Hom}}\nolimits}
\def\Ext{\mathop{\mathrm{Ext}}\nolimits}
\def\SheafHom{\mathop{\mathcal{H}\!\mathit{om}}\nolimits}
\def\SheafExt{\mathop{\mathcal{E}\!\mathit{xt}}\nolimits}
\def\Sch{\mathit{Sch}}
\def\Mor{\mathop{\mathrm{Mor}}\nolimits}
\def\Ob{\mathop{\mathrm{Ob}}\nolimits}
\def\Sh{\mathop{\mathit{Sh}}\nolimits}
\def\NL{\mathop{N\!L}\nolimits}
\def\CH{\mathop{\mathrm{CH}}\nolimits}
\def\proetale{{pro\text{-}\acute{e}tale}}
\def\etale{{\acute{e}tale}}
\def\QCoh{\mathit{QCoh}}
\def\Ker{\mathop{\mathrm{Ker}}}
\def\Im{\mathop{\mathrm{Im}}}
\def\Coker{\mathop{\mathrm{Coker}}}
\def\Coim{\mathop{\mathrm{Coim}}}

% Boxtimes
%
\DeclareMathSymbol{\boxtimes}{\mathbin}{AMSa}{"02}

%
% Macros for moduli stacks/spaces
%
\def\QCohstack{\mathcal{QC}\!\mathit{oh}}
\def\Cohstack{\mathcal{C}\!\mathit{oh}}
\def\Spacesstack{\mathcal{S}\!\mathit{paces}}
\def\Quotfunctor{\mathrm{Quot}}
\def\Hilbfunctor{\mathrm{Hilb}}
\def\Curvesstack{\mathcal{C}\!\mathit{urves}}
\def\Polarizedstack{\mathcal{P}\!\mathit{olarized}}
\def\Complexesstack{\mathcal{C}\!\mathit{omplexes}}
% \Pic is the operator that assigns to X its picard group, usage \Pic(X)
% \Picardstack_{X/B} denotes the Picard stack of X over B
% \Picardfunctor_{X/B} denotes the Picard functor of X over B
\def\Pic{\mathop{\mathrm{Pic}}\nolimits}
\def\Picardstack{\mathcal{P}\!\mathit{ic}}
\def\Picardfunctor{\mathrm{Pic}}
\def\Deformationcategory{\mathcal{D}\!\mathit{ef}}
\usepackage{cleveref}
\crefname{originallemma}{original lemma}{original lemmas}
\Crefname{originallemma}{Original lemma}{Original lemmas}
\makeatletter
\newlabel{derived-lemma-leray-acyclicity}{{015E}{source}{}{sourcepointer-1}{}}
\newlabel{derived-lemma-leray-acyclicity@cref}{{[originallemma][0][]015E}{[1][0][]source}{}{}{}}
\makeatother
\crefname{originallemma}{original lemma}{original lemmas}
\Crefname{originallemma}{Original lemma}{Original lemmas}
\makeatletter
\newlabel{lemma-base-change-map-flat-case}{{02N7}{source}{}{sourcepointer-2}{}}
\newlabel{lemma-base-change-map-flat-case@cref}{{[originallemma][0][]02N7}{[1][0][]source}{}{}{}}
\makeatother
\crefname{originallemma}{original lemma}{original lemmas}
\Crefname{originallemma}{Original lemma}{Original lemmas}
\makeatletter
\newlabel{lemma-cohomology-of-closed}{{09V3}{source}{}{sourcepointer-3}{}}
\newlabel{lemma-cohomology-of-closed@cref}{{[originallemma][0][]09V3}{[1][0][]source}{}{}{}}
\makeatother
\crefname{originallemma}{original lemma}{original lemmas}
\Crefname{originallemma}{Original lemma}{Original lemmas}
\makeatletter
\newlabel{topology-lemma-separated}{{0CY2}{source}{}{sourcepointer-4}{}}
\newlabel{topology-lemma-separated@cref}{{[originallemma][0][]0CY2}{[1][0][]source}{}{}{}}
\makeatother
\title{Proper base change for higher direct images of a module on a ringed space, computed on the fibre}
\author{The Stacks Project Authors}
\date{stacks-project:cohomology.tex:local-snapshot-20261011; GNU FDL 1.2 or later}
\begin{document}
\maketitle
\noindent\textit{Editorial note.} All delivered mathematical text is the author's own, copied from the pinned Stacks Project chapter file cohomology.tex; no statement or proof step was rewritten, reordered or removed.\par\medskip
\noindent\textit{Editorial note.} The proof is complete as written and ends with the construction of the neighbourhood V; no step is referred to the reader.\par\medskip
\noindent\textit{Editorial note.} The counted spans carry 5 navigation stubs and no bibliography citation.\par\medskip
\noindent\textit{Editorial note.} The preamble is the chapter's retained minimal preamble (sources/stacks-cohomology/retained-preamble.tex lines 1-81), a verbatim excerpt of the native preamble.\par\medskip
\noindent\textit{Editorial note.} The label lemma-proper-base-change also occurs in etale-cohomology.tex, but there the statement concerns base change for Rf\_* of a complex with torsion cohomology on the etale site of a proper scheme; hypotheses and conclusion are different, so only the label coincides and this card counts its own proposition.\par\medskip
\noindent\textit{Allowed standard prerequisites (editorial).}\begin{itemize}
\item Base change map for higher direct images. Same chapter, Lemma lemma-base-change-map-flat-case, invoked at proof lines 3073-3074. Supplies the canonical map whose bijectivity is proved.
\item A criterion making the restriction of an injective to a closed subset acyclic. Same chapter, Lemma lemma-cohomology-of-closed, invoked at proof lines 3077-3081 and again at 3091-3092. Supplies that the restrictions of the injective resolution to the fibre are acyclic for global sections, and the identification of the resulting complex.
\item Leray acyclicity: a bounded below complex of acyclic objects computes the derived functor. Derived Categories, Lemma derived-lemma-leray-acyclicity, invoked at proof lines 3082-3084. Supplies that the complex of sections over the fibre represents the derived global sections.
\end{itemize}\medskip
\ifdefined\hypertarget\hypertarget{sourcepointer-1}{}\fi
\noindent\textit{Source reference.} Original lemma 015E: sources/stacks-chow/tags.txt, source lines 1505--1505. Pointer outside the counted span, invoked once for context; the load-bearing argument of the counted statement is the author's proof reproduced in full in the proof span.\par\medskip
\ifdefined\hypertarget\hypertarget{sourcepointer-2}{}\fi
\noindent\textit{Source reference.} Original lemma 02N7: sources/stacks-cohomology/pinned-original-source.tex, source lines 2977--2977. Pointer outside the counted span, invoked once for context; the load-bearing argument of the counted statement is the author's proof reproduced in full in the proof span.\par\medskip
\ifdefined\hypertarget\hypertarget{sourcepointer-3}{}\fi
\noindent\textit{Source reference.} Original lemma 09V3: sources/stacks-cohomology/pinned-original-source.tex, source lines 2848--2848. Pointer outside the counted span, invoked once for context; the load-bearing argument of the counted statement is the author's proof reproduced in full in the proof span.\par\medskip
\ifdefined\hypertarget\hypertarget{sourcepointer-4}{}\fi
\noindent\textit{Source reference.} Original lemma 0CY2: sources/stacks-chow/tags.txt, source lines 16227--16227. Pointer outside the counted span, invoked once for context; the load-bearing argument of the counted statement is the author's proof reproduced in full in the proof span.\par\medskip
\begin{lemma}
\label{lemma-proper-base-change}
Let $f : (X, \mathcal{O}_X) \to (Y, \mathcal{O}_Y)$ be a morphism of
ringed spaces. Let $y \in Y$. Assume that
\begin{enumerate}
\item $f$ is closed,
\item $f$ is separated, and
\item $f^{-1}(y)$ is quasi-compact.
\end{enumerate}
Then for $E$ in $D^+(\mathcal{O}_X)$
we have $(Rf_*E)_y = R\Gamma(f^{-1}(y), E|_{f^{-1}(y)})$ in
$D^+(\mathcal{O}_{Y, y})$.
\end{lemma}

\begin{proof}
The base change map of Lemma \ref{lemma-base-change-map-flat-case}
gives a canonical map $(Rf_*E)_y \to R\Gamma(f^{-1}(y), E|_{f^{-1}(y)})$.
To prove this map is an isomorphism, we represent $E$ by a bounded
below complex of injectives $\mathcal{I}^\bullet$.
Set $Z = f^{-1}(\{y\})$. The assumptions of
Lemma \ref{lemma-cohomology-of-closed}
are satisfied, see Topology, Lemma \ref{topology-lemma-separated}.
Hence the restrictions
$\mathcal{I}^n|_Z$ are acyclic for $\Gamma(Z, -)$.
Thus $R\Gamma(Z, E|_Z)$ is represented by the
complex $\Gamma(Z, \mathcal{I}^\bullet|_Z)$, see
Derived Categories, Lemma \ref{derived-lemma-leray-acyclicity}.
In other words, we have to show the map
$$
\colim_V \mathcal{I}^\bullet(f^{-1}(V))
\longrightarrow
\Gamma(Z, \mathcal{I}^\bullet|_Z)
$$
is an isomorphism. Using Lemma \ref{lemma-cohomology-of-closed}
we see that it suffices to show that the collection of open neighbourhoods
$f^{-1}(V)$ of $Z = f^{-1}(\{y\})$
is cofinal in the system of all open neighbourhoods.
If $f^{-1}(\{y\}) \subset U$ is an open neighbourhood, then as $f$ is closed
the set $V = Y \setminus f(X \setminus U)$ is an open neighbourhood
of $y$ with $f^{-1}(V) \subset U$. This proves the lemma.
\end{proof}

\end{document}
